\section{Explanation}
\label{sec:explanation}

Variant E attaches a two-layer explanation to D's decisions.
The explanation is read-only: E's decisions are D's by construction, and explanation never changes a selection.
% src: A4 (E's decisions equal D's by construction, asserted), X12

\subsection{Layer 1: an exact decomposition of the score}
\label{sec:layer1}
Because the selection score is linear in three terms, the winner-versus-runner-up margin splits exactly into a trust term, a latency term and a load term, and the explanation states which terms decided the choice and why the runner-up lost.
Layer 1 is an identity, not an approximation; it explains the \emph{score}, not why a node's trust has its value.
% src: A12; X15
We check it as bit-for-bit equality with the logged terms and scores, and require that the selected node is the rule's argmax.
On all 6873 decisions of each of A, C and D in the demo scenario (20,619 decisions) there were 0 failures.
% src: A12; F23 (demo scenario, seed 100, draw 7: the exactness identity check is allowed on this scenario, X16)
Figure~\ref{fig:expl} shows one decision as an operator would see it; it is one illustrative decision from the demo scenario (seed 100, draw 7, a tuning seed), not a result.
% src: Figure 6 (fig6_explanation), F25 merged example; X16: illustration only

\begin{figure*}[t]
\centering
\includegraphics[width=0.88\textwidth]{\figdir/fig6_explanation.pdf}
\caption{One decision as an operator sees it (Variant E, the demo scenario: seed 100, draw 7, a tuning seed). The example illustrates the output format of the two layers; it is not a performance or fidelity result. (Generated by \texttt{scripts/make\_figures.py}.)}
\label{fig:expl}
\end{figure*}

\paragraph*{What decides decisions (descriptive only)}
In the same demo scenario (one run per variant, 6873 decisions each), trust is the decisive term---a term worth at least $0.6$ of the winner-versus-runner-up margin---in 15\% of C's decisions and 31\% of D's, and the winner has lower trust than the runner-up in 12.5\% of C's decisions and 18.9\% of D's; load and latency decide most of the rest.
% src: A22; F23 (12.5% = 857/6873, 18.9% = 1296/6873)
This describes how the fixed rule behaved in this one scenario; it carries no interval, is not a performance result and is not used to compare C with D.
% src: A22 (verbatim limits)

\subsection{Layer 2: attribution on the trust head, and its measured limits}
\label{sec:layer2}
Layer 2 applies GNNExplainer to the trust head only, never to the selection rule, and attributes the trust of the winner and of the runner-up to input units; it is an attribution, not a causal account.
% src: A13 (L11); D54 section 1
\TODO{cite GNNExplainer (not in the verified reference set of RELATED\_WORK.md)}
The fidelity test (D54, fixed before any number) perturbs one (node, feature) unit by $0.3$ toward the other end of its range and asks whether the explained RSU's trust moves by at least $0.05$ (\emph{trust change}) or whether the dispatch decision flips (\emph{decision change}); the \emph{top} unit is compared with a random unit drawn from the lower half of the receptive field, and the gap between the two rates is the result.
% src: D54 sections 2, 3 and 5
The floors are $0.20$ for the trust-change gap and $0.05$ for the decision-change gap, and both are stated assumptions.
The test was run on head D twice, on two different injection draws, and on head C once; each pass keeps its own literal label, and the two passes on D are not pooled.
% src: A13 (two passes, each with its own literal label, not pooled); D54 section 5 (floors); X5

\begin{table*}[t]
\centering
\caption{H4 for heads D and C, in this simulation: per-seed mean rate that perturbing the unit changes trust (or flips the decision), and the gap with its 95\% $t$-interval over 5 seeds, with the literal D54 label. C's rows are an artefact of its random-low pool (units C cannot read), not evidence about C's explanations.}
\label{tab:h4}
\small
\begin{tabular}{@{}lllrrrl@{}}
\toprule
Head & Pass & Criterion & Top & Low & Gap [95\% CI] & Label \\
\midrule
% src: A13; F24 (top/low rates, S8; C rows), F25 (labels), F26 H4 block and F27 (S9 replication)
D & S8 (draw 100) & trust change & 0.205 & 0.045 & $+0.160$ [$+0.043$, $+0.277$] & WEAK \\
D & S8 (draw 100) & decision change & 0.065 & 0.032 & $+0.033$ [$+0.0003$, $+0.065$] & WEAK \\
D & S9 (draw 101) & trust change & 0.192 & 0.035 & $+0.158$ [$+0.031$, $+0.284$] & WEAK \\
D & S9 (draw 101) & decision change & 0.070 & 0.020 & $+0.0500$ [$+0.017$, $+0.083$] & FAITHFUL (at the floor) \\
C & S8 (draw 100) & trust change & 0.580 & 0.000 & $+0.580$ [$+0.441$, $+0.720$] & FAITHFUL (artefact) \\
C & S8 (draw 100) & decision change & 0.163 & 0.000 & $+0.163$ [$+0.126$, $+0.199$] & FAITHFUL (artefact) \\
\bottomrule
\end{tabular}
\end{table*}

Table~\ref{tab:h4} gives both passes on D, with C's first-pass rows beside them.
The trust-change gap is below its floor in both passes (WEAK, WEAK).
The decision-change gap is WEAK in the first pass and, in the second, FAITHFUL by the literal $\ge$ rule \emph{exactly at the $0.05$ floor}; we call that result borderline, and we do not average the two passes into a new test.
% src: A13, R3 (borderline); X5 (no pooling)

Three further measured limits apply to Layer 2.
\emph{The masks are unstable}: agreement of the top-attributed unit across three explainer runs is $0.22$ for D, against a pre-registered threshold of $0.80$, so every explanation is one draw and is reported with its run seed.
% src: A13 (top-unit agreement 0.22 vs 0.80); F24, F25
\emph{The neighbour clause is conditional and rare}: the explanation names neighbouring RSUs only when same-segment neighbours carry at least $30\%$ of the attribution, which happened in 10\% of D's explanations (16\% when the explained RSU is degraded).
% src: A13 (10%, 16%); F24 (10.2%, 16.3%)
\emph{It is asynchronous only}: Layer 2 takes 307--598\,ms per decision in the first pass, and 310\,ms at one thread and 675\,ms at ten threads in the replication, against a stated assumption of $100$\,ms (Sec.~\ref{sec:overhead}).
% src: A13 (307-598 ms S8; 310 / 675 ms S9; 100 ms assumption)
The node-local head C also receives a higher fidelity verdict in this test; that is an artefact of its random-low pool, which consists of units C cannot read, and we do not offer it as evidence that C's explanations are better.
% src: A13 (C's higher verdict is an artefact), F25 item 3
